% ====================================================================
% ФАЙЛ: tasks/05_quantum_physics/photoelectric_effect.tex
% ТЕМА: ФОТОЭФФЕКТ. УРАВНЕНИЕ ЭЙНШТЕЙНА ДЛЯ ФОТОЭФФЕКТА
% ====================================================================

% === ЗАДАЧА 1 ===
% Тег: #КВАНТОВАЯ_ФИЗИКА #ФОТОЭФФЕКТ #ЗАПИРАЮЩЕЕ_НАПРЯЖЕНИЕ #КИМ_17
\begin{taskbasic}[code={\label{task:photo_uv_to_green_u_and_work}}]
При исследовании зависимости максимальной кинетической энергии фотоэлектронов от длины волны падающего света фотоэлемент освещали через различные светофильтры. В первой серии опытов использовали светофильтр, пропускающий только ультрафиолетовое излучение, а во второй --- пропускающий только зелёный свет. В каждом опыте наблюдали явление фотоэффекта и измеряли запирающее напряжение.

Как изменились модуль запирающего напряжения и работа выхода фотоэлектронов с поверхности металла при переходе от первой серии опытов ко второй? Для каждой величины определите соответствующий характер её изменения:
1) увеличилась
2) уменьшилась
3) не изменилась

Запишите в таблицу выбранные цифры для каждой физической величины. Цифры в ответе могут повторяться.

\nopagebreak\smallskip
\begin{center}
\small
\begin{tabular}{|c|c|}
\hline
Модуль запирающего напряжения & Работа выхода фотоэлектронов с поверхности металла \\ \hline
 & \\ \hline
\end{tabular}
\end{center}

\kim{17}
\end{taskbasic}
\answer{23}

% === ЗАДАЧА 2 ===
% Тег: #КВАНТОВАЯ_ФИЗИКА #ФОТОЭФФЕКТ #КИНЕТИЧЕСКАЯ_ЭНЕРГИЯ #КИМ_17
\begin{taskbasic}[code={\label{task:photo_green_to_uv_freq_and_ek}}]
При исследовании зависимости максимальной кинетической энергии фотоэлектронов от длины волны падающего света фотоэлемент освещали через различные светофильтры. В первой серии опытов использовали светофильтр, пропускающий только зелёный свет, а во второй --- пропускающий только ультрафиолетовое излучение. В каждом опыте наблюдали явление фотоэффекта и измеряли запирающее напряжение.

Как изменились частота световой волны, падающей на фотоэлемент, и максимальная кинетическая энергия фотоэлектронов при переходе от первой серии опытов ко второй? Для каждой величины определите соответствующий характер её изменения:
1) увеличилась
2) уменьшилась
3) не изменилась

\nopagebreak\smallskip
\begin{center}
\small
\begin{tabular}{|c|c|}
\hline
Частота световой волны, падающей на фотоэлемент & Максимальная кинетическая энергия фотоэлектронов \\ \hline
 & \\ \hline
\end{tabular}
\end{center}

\kim{17}
\end{taskbasic}
\answer{11}

% === ЗАДАЧА 3 ===
% Тег: #КВАНТОВАЯ_ФИЗИКА #ФОТОЭФФЕКТ #ДЛИНА_ВОЛНЫ #КИМ_17
\begin{taskbasic}[code={\label{task:photo_blue_to_green_lambda_and_u}}]
При исследовании зависимости максимальной кинетической энергии фотоэлектронов от частоты падающего света фотоэлемент освещался через светофильтры. В первой серии опытов использовался светофильтр, пропускающий только синий свет, а во второй --- только зелёный. В каждом опыте наблюдали явление фотоэффекта и измеряли запирающее напряжение.

Как изменились длина волны света, падающего на фотоэлемент, и модуль запирающего напряжения при переходе от первой серии опытов ко второй? Для каждой величины определите соответствующий характер её изменения:
1) увеличилась
2) уменьшилась
3) не изменилась

\nopagebreak\smallskip
\begin{center}
\small
\begin{tabular}{|c|c|}
\hline
Длина волны света, падающего на фотоэлемент & Модуль запирающего напряжения \\ \hline
 & \\ \hline
\end{tabular}
\end{center}

\kim{17}
\end{taskbasic}
\answer{12}

% === ЗАДАЧА 4 ===
% Тег: #КВАНТОВАЯ_ФИЗИКА #ФОТОЭФФЕКТ #ЧАСТОТА #КИМ_17
\begin{taskbasic}[code={\label{task:photo_orange_to_blue_freq_and_ek}}]
При исследовании зависимости максимальной кинетической энергии фотоэлектронов от длины волны падающего света фотоэлемент освещали через различные светофильтры. В первой серии опытов использовали светофильтр, пропускающий только оранжевый свет, а во второй --- пропускающий только синий свет. В каждом опыте наблюдали явление фотоэффекта и измеряли запирающее напряжение.

Как изменились частота света, падающего на фотоэлемент, и максимальная кинетическая энергия фотоэлектронов при переходе от первой серии опытов ко второй? Для каждой величины определите соответствующий характер её изменения:
1) увеличилась
2) уменьшилась
3) не изменилась

\nopagebreak\smallskip
\begin{center}
\small
\begin{tabular}{|c|c|}
\hline
Частота света, падающего на фотоэлемент & Максимальная кинетическая энергия фотоэлектронов \\ \hline
 & \\ \hline
\end{tabular}
\end{center}

\kim{17}
\end{taskbasic}
\answer{11}

% === ЗАДАЧА 5 ===
% Тег: #КВАНТОВАЯ_ФИЗИКА #ФОТОЭФФЕКТ #ЭНЕРГИЯ_ФОТОНА #КИМ_17
\begin{taskbasic}[code={\label{task:photo_freq_inc_ephoton_and_work}}]
Во время лабораторной работы ученики изучают зависимость максимальной кинетической энергии фотоэлектронов, вылетающих с фотокатода, от частоты падающего света. В опытах наблюдается явление фотоэффекта.

Частоту падающего света немного увеличивают. Как при этом изменяются энергия фотонов падающего света и работа выхода электронов из материала фотокатода?
1) увеличивается
2) уменьшается
3) не меняется

\nopagebreak\smallskip
\begin{center}
\small
\begin{tabular}{|c|c|}
\hline
Энергия фотонов падающего света & Работа выхода электронов \\ \hline
 & \\ \hline
\end{tabular}
\end{center}

\kim{17}
\end{taskbasic}
\answer{13}

% === ЗАДАЧА 6 ===
% Тег: #КВАНТОВАЯ_ФИЗИКА #ФОТОЭФФЕКТ #ДЛИНА_ВОЛНЫ_УМЕНЬШЕНИЕ #КИМ_17
\begin{taskbasic}[code={\label{task:photo_lambda_dec_ephoton_and_u}}]
Во время лабораторной работы ученики изучают зависимость максимальной кинетической энергии фотоэлектронов, вылетающих с фотокатода, от длины волны падающего света. В опытах наблюдается явление фотоэффекта.

Длину волны падающего света немного уменьшают. Как при этом изменяются энергия фотонов падающего света и модуль запирающего напряжения?
1) увеличивается
2) уменьшается
3) не меняется

\nopagebreak\smallskip
\begin{center}
\small
\begin{tabular}{|c|c|}
\hline
Энергия фотонов падающего света & Модуль запирающего напряжения \\ \hline
 & \\ \hline
\end{tabular}
\end{center}

\kim{17}
\end{taskbasic}
\answer{11}

% === ЗАДАЧА 7 ===
% Тег: #КВАНТОВАЯ_ФИЗИКА #ФОТОЭФФЕКТ #ЧАСТОТА_УВЕЛИЧЕНИЕ #КИМ_17
\begin{taskbasic}[code={\label{task:photo_freq_inc_ek_and_work}}]
В лабораторной работе ученик изучает зависимость максимальной кинетической энергии фотоэлектронов, вылетающих с фотокатода, от частоты падающего света. В опытах наблюдается явление фотоэффекта.

Частоту падающего света в опыте немного увеличивают. Как при этом изменяются максимальная кинетическая энергия фотоэлектронов и работа выхода фотоэлектронов из металла фотокатода?
1) увеличивается
2) уменьшается
3) не изменяется

\nopagebreak\smallskip
\begin{center}
\small
\begin{tabular}{|c|c|}
\hline
Максимальная кинетическая энергия фотоэлектронов & Работа выхода фотоэлектронов \\ \hline
 & \\ \hline
\end{tabular}
\end{center}

\kim{17}
\end{taskbasic}
\answer{13}

% === ЗАДАЧА 8 ===
% Тег: #КВАНТОВАЯ_ФИЗИКА #ФОТОЭФФЕКТ #ИНТЕНСИВНОСТЬ #КИМ_17
\begin{taskbasic}[code={\label{task:photo_intensity_inc_u_and_redline}}]
В лабораторной работе ученик изучает зависимость максимальной кинетической энергии фотоэлектронов, вылетающих с фотокатода, от частоты падающего света. В опытах наблюдается явление фотоэффекта.

Интенсивность падающего света в опыте немного увеличивают, не изменяя его частоты. Как при этом изменяются модуль запирающего напряжения и длина волны, соответствующая «красной границе» фотоэффекта?
1) увеличивается
2) уменьшается
3) не изменяется

\nopagebreak\smallskip
\begin{center}
\small
\begin{tabular}{|c|c|}
\hline
Модуль запирающего напряжения & Длина волны, соответствующая «красной границе» фотоэффекта \\ \hline
 & \\ \hline
\end{tabular}
\end{center}

\kim{17}
\end{taskbasic}
\answer{33}

% === ЗАДАЧА 9 ===
% Тег: #КВАНТОВАЯ_ФИЗИКА #ФОТОЭФФЕКТ #СИНИЙ_ЗЕЛЕНЫЙ_FREQ_U #КИМ_17
\begin{taskbasic}[code={\label{task:photo_blue_to_green_freq_and_u}}]
При исследовании зависимости кинетической энергии фотоэлектронов от частоты падающего света фотоэлемент освещался через светофильтры. В первой серии опытов использовался светофильтр, пропускающий только синий свет, а во второй --- только зелёный. В каждом опыте наблюдали явление фотоэффекта и измеряли запирающее напряжение.

Как изменятся частота световой волны и модуль запирающего напряжения при переходе от первой серии опытов ко второй? Для каждой величины определите соответствующий характер её изменения:
1) увеличится
2) уменьшится
3) не изменится

\nopagebreak\smallskip
\begin{center}
\small
\begin{tabular}{|c|c|}
\hline
Частота волны света, падающего на фотоэлемент & Модуль запирающего напряжения \\ \hline
 & \\ \hline
\end{tabular}
\end{center}

\kim{17}
\end{taskbasic}
\answer{22}

% === ЗАДАЧА 10 ===
% Тег: #КВАНТОВАЯ_ФИЗИКА #ФОТОЭФФЕКТ #СИНИЙ_ЗЕЛЕНЫЙ_LAMBDA_WORK #КИМ_17
\begin{taskbasic}[code={\label{task:photo_blue_to_green_lambda_and_work}}]
При исследовании зависимости кинетической энергии фотоэлектронов от частоты падающего света фотоэлемент освещался через светофильтры. В первой серии опытов использовался светофильтр, пропускающий только синий свет, а во второй --- только зелёный. В каждом опыте наблюдали явление фотоэффекта и измеряли запирающее напряжение.

Как изменятся длина световой волны и работа выхода электронов из металла при переходе от первой серии опытов ко второй? Для каждой величины определите соответствующий характер её изменения:
1) увеличится
2) уменьшится
3) не изменится

\nopagebreak\smallskip
\begin{center}
\small
\begin{tabular}{|c|c|}
\hline
Длина волны света, падающего на фотоэлемент & Работа выхода \\ \hline
 & \\ \hline
\end{tabular}
\end{center}

\kim{17}
\end{taskbasic}
\answer{13}

% ====================================================================
% ФАЙЛ: tasks/05_quantum_physics/photoelectric_effect_part2.tex
% ТЕМА: ФОТОЭФФЕКТ. ГРАФИКИ И СВОЙСТВА ФОТОЭЛЕМЕНТА
% ====================================================================

% === ЗАДАЧА 1 ===
% Тег: #КВАНТОВАЯ_ФИЗИКА #ФОТОЭФФЕКТ #ГРАФИКИ #КИМ_17
\begin{taskbasic}[code={\label{task:photo_graphs_ek_nu_e_lambda}}]
Установите соответствие между графиками, представленными на рисунках, и зависимостями, которые они могут выражать.

К каждой позиции первого столбца подберите соответствующую позицию из второго столбца и запишите в таблицу выбранные цифры под соответствующими буквами.

\nopagebreak\smallskip
\begin{center}
\begin{tikzpicture}[scale=0.9, every node/.style={font=\footnotesize}]
    \draw[xstep=1, ystep=1, gray!30, thin] (0,0) grid (3,2);
    
    % График А
    \draw[-{Stealth[scale=0.8]}, thick, black] (0,0) -- (0,2.3);
    \draw[-{Stealth[scale=0.8]}, thick, black] (0,0) -- (3.3,0);
    \draw[ultra thick, brandPrimary] (1.0,0) -- (2.8,1.2);
    \node[below left] at (0,0) {0};
    \node[left] at (-0.2,1.8) {\textbf{А)}};
\end{tikzpicture}
\qquad\qquad
\begin{tikzpicture}[scale=0.9, every node/.style={font=\footnotesize}]
    \draw[xstep=1, ystep=1, gray!30, thin] (0,0) grid (3,2);
    
    % График Б
    \draw[-{Stealth[scale=0.8]}, thick, black] (0,0) -- (0,2.3);
    \draw[-{Stealth[scale=0.8]}, thick, black] (0,0) -- (3.3,0);
    \draw[ultra thick, brandPrimary] plot[domain=0.4:3.0, samples=100] (\x, {0.8/\x});
    
    \draw[dashed, gray] (0.5,0) -- (0.5,1.6) -- (0,1.6);
    \draw[dashed, gray] (1.0,0) -- (1.0,0.8) -- (0,0.8);
    \draw[dashed, gray] (2.0,0) -- (2.0,0.4) -- (0,0.4);
    
    \node[below left] at (0,0) {0};
    \node[left] at (-0.2,1.8) {\textbf{Б)}};
\end{tikzpicture}
\end{center}

\smallskip
\begin{tabular}{p{0.35\textwidth}p{0.6\textwidth}}
\textbf{ГРАФИКИ} & \textbf{ЗАВИСИМОСТИ} \\
А) График А & 1) зависимость энергии фотона от длины волны \\
Б) График Б & 2) зависимость максимальной энергии фотоэлектронов от частоты света \\
& 3) зависимость энергии фотона от частоты света \\
& 4) зависимость силы фототока от напряжения между электродами при неизменной освещённости
\end{tabular}

\kim{17}
\end{taskbasic}
\answer{21}

% === ЗАДАЧА 2 ===
% Тег: #КВАНТОВАЯ_ФИЗИКА #ФОТОЭФФЕКТ #ГРАФИКИ #КИМ_17
\begin{taskbasic}[code={\label{task:photo_graphs_e_nu_i_u}}]
Установите соответствие между графиками, представленными на рисунках, и зависимостями, которые они могут выражать.

К каждой позиции первого столбца подберите соответствующую позицию из второго столбца и запишите в таблицу выбранные цифры под соответствующими буквами.

\nopagebreak\smallskip
\begin{center}
\begin{tikzpicture}[scale=0.9, every node/.style={font=\footnotesize}]
    \draw[xstep=1, ystep=1, gray!30, thin] (0,0) grid (3,2);
    
    % График А
    \draw[-{Stealth[scale=0.8]}, thick, black] (0,0) -- (0,2.3);
    \draw[-{Stealth[scale=0.8]}, thick, black] (0,0) -- (3.3,0);
    \draw[ultra thick, brandPrimary] (0,0) -- (2.8,1.6);
    \node[below left] at (0,0) {0};
    \node[left] at (-0.2,1.8) {\textbf{А)}};
\end{tikzpicture}
\qquad\qquad
\begin{tikzpicture}[scale=0.9, every node/.style={font=\footnotesize}]
    \draw[xstep=1, ystep=1, gray!30, thin] (0,0) grid (3,2);
    
    % График Б (ВАХ)
    \draw[-{Stealth[scale=0.8]}, thick, black] (0,0) -- (0,2.3);
    \draw[-{Stealth[scale=0.8]}, thick, black] (0,0) -- (3.3,0);
    \draw[ultra thick, brandPrimary] plot[domain=0:2.8, samples=100] (\x, {1.6/(1 + exp(-3*(\x-1.0)))});
    \draw[dashed, gray] (0,1.6) -- (2.8,1.6);
    \node[below left] at (0,0) {0};
    \node[left] at (-0.2,1.8) {\textbf{Б)}};
\end{tikzpicture}
\end{center}

\smallskip
\begin{tabular}{p{0.35\textwidth}p{0.6\textwidth}}
\textbf{ГРАФИКИ} & \textbf{ЗАВИСИМОСТИ} \\
А) График А & 1) зависимость энергии фотона от его длины волны \\
Б) График Б & 2) зависимость максимальной энергии фотоэлектронов от частоты света \\
& 3) зависимость энергии фотона от частоты света \\
& 4) зависимость силы фототока от напряжения между электродами при неизменной освещённости
\end{tabular}

\kim{17}
\end{taskbasic}
\answer{34}

% === ЗАДАЧА 3 ===
% Тег: #КВАНТОВАЯ_ФИЗИКА #ФОТОЭФФЕКТ #ЗЕЛЕНЫЙ_СИНИЙ_LAMBDA_REDLINE #КИМ_17
\begin{taskbasic}[code={\label{task:photo_green_to_blue_lambda_and_redline}}]
При исследовании зависимости кинетической энергии фотоэлектронов от частоты падающего света фотоэлемент освещался через светофильтры. В первой серии опытов использовался светофильтр, пропускающий только зелёный свет, а во второй --- пропускающий только синий. В каждом опыте наблюдали явление фотоэффекта.

Как изменились длина волны падающего на фотоэлемент света и «красная граница» фотоэффекта при переходе от первой серии опытов ко второй?
1) увеличилась
2) уменьшилась
3) не изменилась

\nopagebreak\smallskip
\begin{center}
\small
\begin{tabular}{|c|c|}
\hline
Длина волны падающего на фотоэлемент света & «Красная граница» фотоэффекта \\ \hline
 & \\ \hline
\end{tabular}
\end{center}

\kim{17}
\end{taskbasic}
\answer{23}

% === ЗАДАЧА 4 ===
% Тег: #КВАНТОВАЯ_ФИЗИКА #ФОТОЭФФЕКТ #СИНИЙ_ЖЕЛТЫЙ_FREQ_WORK #КИМ_17
\begin{taskbasic}[code={\label{task:photo_blue_to_yellow_freq_and_work}}]
При исследовании зависимости кинетической энергии фотоэлектронов от частоты падающего света фотоэлемент освещался через светофильтры. В первой серии опытов использовался светофильтр, пропускающий только синий свет, а во второй --- пропускающий только жёлтый. В каждом опыте наблюдали явление фотоэффекта.

Как изменились частота света, падающего на фотоэлемент, и работа выхода электронов при переходе от первой серии опытов ко второй?
1) увеличилась
2) уменьшилась
3) не изменилась

\nopagebreak\smallskip
\begin{center}
\small
\begin{tabular}{|c|c|}
\hline
Частота света, падающего на фотоэлемент & Работа выхода электронов \\ \hline
 & \\ \hline
\end{tabular}
\end{center}

\kim{17}
\end{taskbasic}
\answer{23}

% === ЗАДАЧА 5 ===
% Тег: #КВАНТОВАЯ_ФИЗИКА #ФОТОЭФФЕКТ #U_INC_LAMBDA_WORK #КИМ_17
\begin{taskbasic}[code={\label{task:photo_u_inc_lambda_and_work}}]
Монохроматический свет с длиной волны $\lambda$ падает на поверхность металла, вызывая фотоэффект. После изменения энергии падающих фотонов модуль запирающего напряжения $U_{\text{зап}}$ увеличился. Как изменились при этом длина волны $\lambda$ падающего света и работа выхода фотоэлектронов с поверхности металла?

Для каждой величины определите соответствующий характер изменения:
1) увеличилась
2) уменьшилась
3) не изменилась

\nopagebreak\smallskip
\begin{center}
\small
\begin{tabular}{|c|c|}
\hline
Длина волны $\lambda$ падающего света & Работа выхода фотоэлектронов \\ \hline
 & \\ \hline
\end{tabular}
\end{center}

\kim{17}
\end{taskbasic}
\answer{23}

% === ЗАДАЧА 6 ===
% Тег: #КВАНТОВАЯ_ФИЗИКА #ФОТОЭФФЕКТ #U_DEC_LAMBDA_EK #КИМ_17
\begin{taskbasic}[code={\label{task:photo_u_dec_lambda_and_ek}}]
Монохроматический свет с длиной волны $\lambda$ падает на поверхность металла, вызывая фотоэффект. После изменения энергии падающих фотонов модуль запирающего напряжения $U_{\text{зап}}$ уменьшился. Как изменились при этом длина волны $\lambda$ падающего света и максимальная кинетическая энергия фотоэлектронов?

Для каждой величины определите соответствующий характер изменения:
1) увеличилась
2) уменьшилась
3) не изменилась

\nopagebreak\smallskip
\begin{center}
\small
\begin{tabular}{|c|c|}
\hline
Длина волны $\lambda$ падающего света & Максимальная кинетическая энергия фотоэлектронов \\ \hline
 & \\ \hline
\end{tabular}
\end{center}

\kim{17}
\end{taskbasic}
\answer{12}